% Part: first-order-logic
% Chapter: tableaux
% Section: provability-quantifiers

\documentclass[../../../include/open-logic-section]{subfiles}

\begin{document}

\olfileid{fol}{tab}{qpr}
\olsection{\usetoken{S}{derivability} en de kwantoren (‘voor alle’ en ‘er bestaat’)}

\begin{explain}
  Voor het bewijs van de volledigheidsstelling moeten de tableauregels
  ook de feiten over~$\Proves$ opleveren die we hier bewijzen.
\end{explain}

\begin{thm}
\ollabel{thm:strong-generalization} Stel dat $c$ een constante is die
niet voorkomt in $\Gamma$ en ook niet in~$!A(x)$. Als
$\Gamma \Proves !A(c)$, dan geldt
$\Gamma \Proves \lforall[x][!A(x)]$.
\end{thm}

\begin{proof}
Neem aan dat $\Gamma \Proves !A(c)$. Dat betekent dat er
$!B_1$, \dots, $!B_n \in \Gamma$ zijn en dat er een gesloten
!!{tableau} is voor
\begin{align*}
\{ \sFmla{\False}{!A(c)}, &
\sFmla{\True}{!B_1}, \dots, \sFmla{\True}{!B_n} \}.
\intertext{We moeten ook een gesloten !!{tableau} maken voor}
\{ \sFmla{\False}{\lforall[x][!A(x)]}, &
\sFmla{\True}{!B_1}, \dots, \sFmla{\True}{!B_n} \}.
\end{align*}
Verander het gesloten !!{tableau} als volgt. Vervang de eerste aanname
door $\sFmla{\False}{\lforall[x][!A(x)]}$. Voeg na de aannamen
$\sFmla{\False}{!A(c)}$ in.
\begin{center}
\begin{tableau}{not line numbering}
  [\sFmla{\False}{\formula{A}(c)}
    [\sFmla{\True}{\formula{B}_1}
      [\vdots
      [\sFmla{\True}{\formula{B}_n}, 
        [][][]]]]]
\end{tableau}\qquad
\begin{tableau}{not line numbering}
  [\sFmla{\False}{\lforall[x][\formula{A}(x)]}
    [\sFmla{\True}{\formula{B}_1}
      [\vdots
        [\sFmla{\True}{\formula{B}_n},
          [\sFmla{\False}{\formula{A}(c)}
        [][][]]]]]]
\end{tableau}
\end{center}
Waarom blijft het tableau gesloten? Alle !!{sentence}s die eerst als
aannamen beschikbaar waren, staan nog steeds bovenaan het
!!{tableau}. Ook de nieuwe regel is juist. Hij volgt door
$\TRule{\False}{\lforall}$ toe te passen. De !!{constant}~$c$ komt
namelijk niet voor in $!B_1$, \dots, $!B_n$ of
$\lforall[x][!A(x)]$. Hij komt dus niet voor boven de ingevoegde regel
in het nieuwe !!{tableau}.
\end{proof}

\begin{prop}
\ollabel{prop:provability-quantifiers}
\begin{tagenumerate}{prvEx,prvAll}
\tagitem{prvEx}{$!A(t) \Proves \lexists[x][!A(x)]$.}{}
\tagitem{prvAll}{$\lforall[x][!A(x)] \Proves !A(t)$.}{}
\end{tagenumerate}
\end{prop}

\begin{proof}
\begin{tagenumerate}{prvEx,prvAll}
\tagitem{prvEx}{Dit is een gesloten !!{tableau} voor
  $\sFmla{\False}{\lexists[x][!A(x)]}, \sFmla{\True}{!A(t)}$:
  \begin{oltableau}
    [\sFmla{\False}{\lexists[x][\formula{A}(x)]}, just = \TAss
      [\sFmla{\True}{\formula{A}(t)}, just = \TAss
        [\sFmla{\False}{\formula{A}(t)}, just = {\TRule{\False}{\lexists}[1]}, close]]]
    \end{oltableau}}{}
\tagitem{prvAll}{Dit is een gesloten !!{tableau} voor
  $\sFmla{\False}{\formula{A}(t)}, \sFmla{\True}{\lforall[x][!A(x)]}, $:
  \begin{oltableau}
    [\sFmla{\False}{\formula{A}(t)}, just = \TAss
      [\sFmla{\True}{\lforall[x][\formula{A}(x)]}, just = \TAss
        [\sFmla{\True}{\formula{A}(t)}, just = {\TRule{\True}{\lforall}[2]}, close]]]
    \end{oltableau}}{}
\end{tagenumerate}
\end{proof}

\end{document}
